\chapter{Discussion} \label{cha:discussion}

The \textbf{Discussion} chapter interprets and contextualizes your results. This is where you explain \textit{what your findings mean}, \textit{why they matter}, and \textit{how they relate} to previous research and the original purpose of the study. You may also reflect on methodological choices, limitations, and possible directions for future work.

\subsection*{Goals of the Chapter}

The discussion should:
\begin{itemize}
  \item Answer the research questions based on your results
  \item Reflect critically on the strengths and weaknesses of your study
  \item Show how your findings relate to the background and literature
  \item Consider the broader implications of your work
\end{itemize}

\subsection*{Suggested Structure}

\begin{enumerate}
  \item \textbf{Summary of Main Results} \\
  Start by briefly restating your most important results in relation to the research questions. Focus on patterns, contrasts, and significant observations.

  \item \textbf{Interpretation and Comparison} \\
  Analyze what the results suggest. Discuss:
  \begin{itemize}
    \item Why certain models or methods performed as they did
    \item Surprising or unexpected outcomes
    \item How your findings align with or diverge from previous studies
    \item Possible explanations grounded in theory, language typology, model architecture, or dataset characteristics
  \end{itemize}

  \item \textbf{Methodological Reflections} \\
  Reflect on your research design:
  \begin{itemize}
    \item Was the data sufficient and appropriate?
    \item Were the tools and metrics adequate?
    \item Could another approach have led to different or better results?
    \item Were there any sources of bias, simplification, or limitations?
  \end{itemize}
  Be honest and precise—critical reflection is a strength, not a weakness.

  \item \textbf{Limitations} \\
  Identify the key limitations of your study:
  \begin{itemize}
    \item Scope limitations (e.g., only one language family or corpus)
    \item Data constraints (e.g., size, quality, imbalance)
    \item Technical limitations (e.g., computational resources, model support)
  \end{itemize}
  Mention how these might affect the validity or generalizability of your findings.

  \item \textbf{Implications and Applications} \\
  Discuss the broader relevance of your findings:
  \begin{itemize}
    \item For future research in the area of study, e.g. for NLP, speech technology, neuroscience, or linguistics
    \item For technology development (e.g., tools for under-resourced languages)
    \item For societal or ethical considerations (e.g., inclusion, fairness)
  \end{itemize}

  \item \textbf{Suggestions for Future Work} \\
  Suggest logical next steps that could build on your work:
  \begin{itemize}
    \item Replication studies on new datasets or languages
    \item Extension of methods to related tasks
    \item Addressing limitations you identified
    \item Scaling or adapting the approach for real-world applications
  \end{itemize}
\end{enumerate}

\subsection*{Best Practices}

\begin{itemize}
  \item Use headings or subheadings to structure your discussion clearly.
  \item Be critical but constructive—highlight what worked and what didn’t.
  \item Avoid repeating raw results—focus instead on interpretation and reflection.
  \item Keep your claims proportionate to your evidence—don’t overgeneralize.
  \item Refer back to your literature review where appropriate to show connections.
\end{itemize}
